% !TITLE 将进酒
% !AUTHOR 李白
% !DYNASTY 唐
% !ORDER 11

\begin{poem}
君不见黄河之水天上来，奔流到海不复回。
君不见高堂明镜悲白发，朝如青丝暮成雪。
人生得意须尽欢，莫使金樽空对月。
天生我材必有用，千金散尽还复来。
烹羊宰牛且为乐，会须一饮三百杯。

岑夫子，丹丘生，将进酒，杯莫停。
与君歌一曲，请君为我倾耳听。
钟鼓馔玉不足贵，但愿长醉不复醒。
古来圣贤皆寂寞，惟有饮者留其名。
陈王昔时宴平乐，斗酒十千恣欢谑。
主人何为言少钱，径须沽取对君酌。
五花马，千金裘，呼儿将出换美酒，与尔同销万古愁。
\end{poem}

\begin{appreciation}
这是李白最豪放的诗，也是他内心最痛苦的诗。写这首诗时，他已经离开长安，在梁宋一带漫游。表面上是在劝酒，骨子里是在与绝望搏斗。

君不见黄河之水天上来，奔流到海不复回——黄河从天上流下来，奔流到海再也不回头。这是时间的比喻：时间像黄河一样，一旦流逝就永不复返。君不见高堂明镜悲白发，朝如青丝暮成雪——早上还是黑发，晚上就变成了白雪。这种夸张不是修辞，是一个中年人对时间流逝的真实恐惧。李白写这首诗时大约五十岁，在古代已算老年。他终于意识到了生命的有限。

人生得意须尽欢，莫使金樽空对月——既然生命短暂，那就及时行乐。天生我材必有用，千金散尽还复来——这两句是全诗最著名的豪言。但请注意，它出现在一个极度悲观的语境中。正是因为黄河不回、白发如雪，才需要用天生我材必有用来安慰自己。这种在绝望中硬撑起来的自信，比一帆风顺时的豪言更令人动容。

钟鼓馔玉不足贵，但愿长醉不复醒——钟鸣鼎食的富贵生活不值得珍贵，我只想长醉不醒。古来圣贤皆寂寞，惟有饮者留其名——古代的圣贤都很寂寞，只有饮酒的人留下了名字。这不是真的在赞美饮酒，而是在说：既然清醒是如此痛苦，不如醉去。李白的醉，很多时候是一种自我保护——当现实太残酷时，酒是最好的避风港。

五花马，千金裘，呼儿将出换美酒，与尔同销万古愁——最后一句点出了全诗的真意。万古愁，三个字把个人的忧愁扩展到了人类共同的命运。李白不是在为自己发愁，他是在为生命的短暂、为理想的无着、为古往今来所有不被理解的人发愁。这种愁，不是一杯酒可以销掉的，但他还是要销——因为不销，人就活不下去了。
\end{appreciation}
